\centerline{\bf \Huge Степень вхождения}

\begin{flushright}
        \scriptsize{Я чувствую себя пакетиком чая, который решили заварить\\ Аква из КоноСуба}
\end{flushright}

Пусть $p$ - простое число. Каждое целое ненулевое число $n$ можно единственным образом представить в виде $n=p^k m$, где $m$ не делится на $p$ (возможно, $k=0$ ). Число $k$ называется степенью вхождения простого $p$ в $n$. Обозначение: $k=\nu_p(n)$.

\textbf{\large Свойства:}

(a) $\nu_p(a b)=\nu_p(a)+\nu_p(b)$;

(б) $\nu_p(a \pm b) \geqslant \min \left\{\nu_p(a), \nu_p(b)\right\}$, причём если $\nu_p(a) \neq \nu_p(b)$, то достигается равенство.

\problem{0.1}
Натуральные числа $a, b, c, d$ таковы, что $a b c, a b d, a c d, b c d$ - точные квадраты. Докажите, что $a, b, c, d$ - тоже точные квадраты.

\problem{0.2}
Натуральные числа $a$ и $b$ таковы, что сумма $\frac{a^2}{b}+\frac{b^2}{a}$ целая. Докажите, что оба слагаемых целые.

\problem{1}
Степень вхождения простого числа $p$ в натуральное $a$ равна $\alpha$, а степень вхождения $p$ в натуральное $b$ равна $\beta$. Докажите, что степень вхождения $p$ в НОД $(a, b)$ равна $\min (\alpha, \beta)$, а степень вхождения $p$ в $\operatorname{HOK}(a, b)$ равна $\max (\alpha, \beta)$.

\problem{2}
(a) \textbf{Формула Лежандра.} Докажите, что степень вхождения простого числа $p$ в $n!$ может быть вычислена по формуле

$$
\nu_p(n!)=\left[\frac{n}{p}\right]+\left[\frac{n}{p^2}\right]+\left[\frac{n}{p^3}\right]+\ldots
$$

(b)
\textbf{Gasanov's Lemma.}
Let $n$ be a natural number and $p$ be a prime greater than 2 , then $\nu_p(n!) \leq\left[\frac{n}{p-1}\right]$.

(c)
\textbf{Gasanov's Theorem.}
Let $a_1, a_2, \ldots, a_k$ positive integers such that $a_1+a_2+\ldots+a_k=$ $k(p-1)$, where $p$ is prime number and $k>2$, then the product $a_{1}!\cdot a_{2}!\cdot \ldots \cdot a_{k}!$ will be an exact $k$-th power of a natural number if and only if

$$
a_1=a_2=\ldots=a_k=p-1
$$

\begin{flushright}
        \scriptsize{за постулат Бертрана бан!}
\end{flushright}

\newpage

\hspace{3mm}

\problem{3}
Докажите, что число $1+\frac{1}{2}+\frac{1}{3}+\ldots+\frac{1}{n}$ не является целым ни при каком натуральном $n \geqslant 2$.


\begin{flushright}
        \scriptsize{тут тоже за постулат Бертрана бан!}
\end{flushright}

\problem{4}
Существует ли бесконечная арифметическая прогрессия с ненулевой разностью, состоящая только из степеней (выше первой!!!) натуральных чисел?

\problem{5}
Для натурального числа $n$ обозначим через $S_n$ наименьшее общее кратное всех чисел $1,2,3, \ldots, n$. Существует ли такое натуральное число $m$, что $S_{m+1}=4 S_m ?$


\problem{6}
По кругу стоят $10^{1000}$ натуральных чисел. Между каждыми двумя соседними числами записали их наименьшее общее кратное. Могут ли эти наименьшие общие кратные образовать $10^{1000}$ последовательных чисел (расположенных в каком-то порядке)?

\problem{7}
Петя выбрал несколько последовательных натуральных чисел и каждое записал либо красным, либо синим карандашом (оба цвета присутствуют). Может ли сумма наименьшего общего кратного всех красных чисел и наименьшего общего кратного всех синих чисел являться степенью двойки?

\problem{8}
Докажите, что

(a)
$$
\nu_p(C_n^k)>\nu_p(n)-\frac{k}{p-1}
$$

(б)
\textbf{LTE-лемма.} Пусть $x$ и $y$ - различные ненулевые целые числа, $p$ - нечетное простое число, не являющееся делителем $x$ и $y$ и такое, что $p \mid x \pm y$. Тогда для любого натурального $n$ выполнено равенство

$$
\nu_p(x^n \pm y^n) = \nu_p(x \pm y)+\nu_p (n)
$$